\section{Assignments}

\subsection{Assignment 1}

\subsubsection*{Problem 1}

\term{Question.} Consider the following two Monte Carlo schemes for estimating

$\mu=\mathbb{E}[h(X)+f(X)]$.
The total sample size is $2n$ in both schemes. Scheme 1 (same random numbers): generate $2n$ i.i.d. copies $X_1,\ldots,X_{2n}$ and use

$\widehat\mu_1=\frac1{2n}\sum_{i=1}^{2n}[h(X_i)+f(X_i)]$.
Scheme 2 (different random numbers): write $\mu=\mathbb{E}[h(X)]+\mathbb{E}[f(X)]$, generate $X_1,\ldots,X_n,Y_1,\ldots,Y_n$ i.i.d., and use

$\widehat\mu_2=\frac1n\sum_{i=1}^nh(X_i)+\frac1n\sum_{i=1}^nf(Y_i)$.
Show that both estimators are unbiased and
$\operatorname{var}(\widehat\mu_1)-\operatorname{var}(\widehat\mu_2)=-\frac1{2n}\operatorname{var}[h(X)-f(X)]$.

Hence deduce that $\operatorname{var}(\widehat\mu_1)\leq\operatorname{var}(\widehat\mu_2)$.

\term{Solution.}

$\mathbb{E}[\widehat\mu_1]=\frac1{2n}\sum_{i=1}^{2n}\{\mathbb{E}[h(X_i)]+\mathbb{E}[f(X_i)]\}=\mathbb{E}[h(X)]+\mathbb{E}[f(X)]=\mu$,

$\mathbb{E}[\widehat\mu_2]=\frac1n\sum_{i=1}^n\mathbb{E}[h(X_i)]+\frac1n\sum_{i=1}^n\mathbb{E}[f(Y_i)]=\mathbb{E}[h(X)]+\mathbb{E}[f(X)]=\mu$.

Thus both are unbiased. Moreover,

$\operatorname{var}(\widehat\mu_1)=\frac1{4n^2}\sum_{i=1}^{2n}\operatorname{var}[h(X_i)+f(X_i)]$

$=\frac1{2n}\{\operatorname{var}[h(X)]+2\operatorname{cov}[h(X),f(X)]+\operatorname{var}[f(X)]\}$,

whereas independence of the two samples gives

$\operatorname{var}(\widehat\mu_2)=\frac1n\operatorname{var}[h(X)]+\frac1n\operatorname{var}[f(X)]$.

Therefore

$\operatorname{var}(\widehat\mu_1)-\operatorname{var}(\widehat\mu_2)$

$=\frac1{2n}\{-\operatorname{var}[h(X)]+2\operatorname{cov}[h(X),f(X)]-\operatorname{var}[f(X)]\}$

$=-\frac1{2n}\operatorname{var}[h(X)-f(X)]\leq0$.

Hence $\operatorname{var}(\widehat\mu_1)\leq\operatorname{var}(\widehat\mu_2)$.

\subsubsection*{Problem 2}

\term{Question.} Given that $S(t)$ is a geometric Brownian motion satisfying

$\frac{dS(t)}{S(t)}=r\,dt+\sigma\,dW(t),\qquad S(0)=S$,

where $W(t)$ is a standard Brownian motion under the risk-neutral measure $\mathbb{Q}$: (a) let $Y(t)=S^\alpha(t)$ and find its SDE; (b) using that SDE, suggest two Monte Carlo methods to estimate

$\mathbb{E}^{\mathbb Q}\left[e^{-rT}\max\left(\frac1m\sum_{i=1}^mY(t_i)-K,0\right)\right],\qquad 0<t_1<\cdots<t_m<T$.

One method must give a biased estimator and the other an unbiased estimator.

\term{Solution.} (a) By It\^o's lemma,

$dY(t)=\frac{\partial Y}{\partial t}dt+\frac{\partial Y}{\partial S}dS(t)+\frac12\frac{\partial^2Y}{\partial S^2}(dS(t))^2$

$=\alpha S^{\alpha-1}(t)dS(t)+\frac12\alpha(\alpha-1)S^{\alpha-2}(t)(dS(t))^2$

$=\alpha S^\alpha(t)[r\,dt+\sigma\,dW(t)]+\frac12\alpha(\alpha-1)S^\alpha(t)[r\,dt+\sigma\,dW(t)]^2$

$=S^\alpha(t)\left[\left(\alpha r+\frac12\alpha(\alpha-1)\sigma^2\right)dt+\alpha\sigma\,dW(t)\right]$

$=Y(t)\left[\left(\alpha r+\frac12\alpha(\alpha-1)\sigma^2\right)dt+\alpha\sigma\,dW(t)\right]$,
with $Y(0)=S^\alpha$.

(b) By Monte Carlo,
$\widehat V_n=\frac1n\sum_{k=1}^nV_k$,

where
$V_k=\max\left(\frac1m\sum_{i=1}^mY^{(k)}(t_i)-K,0\right)$

and $\{Y^{(k)}(t_1),\ldots,Y^{(k)}(t_m)\}$ is the $k$th path of $Y(t)$.

\term{Method 1 (Biased Estimator).} Discretizing the SDE gives

$Y(t_i)-Y(t_{i-1})\approx\left(\alpha r+\frac12\alpha(\alpha-1)\sigma^2\right)Y(t_{i-1})(t_i-t_{i-1})$

$\qquad+\alpha\sigma Y(t_{i-1})[W(t_i)-W(t_{i-1})]$,

so
$Y(t_i)\approx Y(t_{i-1})+\left(\alpha r+\frac12\alpha(\alpha-1)\sigma^2\right)Y(t_{i-1})(t_i-t_{i-1})$

$\qquad+\alpha\sigma Y(t_{i-1})\sqrt{t_i-t_{i-1}}Z_i,\qquad Z_i\sim N(0,1)$.

Thus the $k$th path is generated by applying this recursion with independent $Z_i^{(k)}$.

\term{Method 2 (Unbiased Estimator).} Solving the SDE gives

$Y(t_i)=Y(t_{i-1})\exp(A_i+B_i)$,

$A_i=\left(\alpha r+\frac12\alpha(\alpha-1)\sigma^2-\frac12\alpha^2\sigma^2\right)(t_i-t_{i-1})$,

$B_i=\alpha\sigma[W(t_i)-W(t_{i-1})]$.

$=Y(t_{i-1})\exp\left[\left(\alpha r-\frac12\alpha\sigma^2\right)(t_i-t_{i-1})+\alpha\sigma\sqrt{t_i-t_{i-1}}Z_i\right]$.

Thus the $k$th path is generated from the exact recursion with independent $Z_i^{(k)}$.

\subsubsection*{Problem 3}

\term{Question.} Simulate $X$ by the inverse transform algorithm

$X=\operatorname{int}(10U(1-U))$,

where $\operatorname{int}(x)$ is the integral part of $x$ and $U\sim\operatorname{Unif}(0,1)$. Find $P(X=0)$, $P(X=1)$ and $P(X=2)$.

\term{Solution.} Since $U\in[0,1]$, $10U(1-U)\in[0,2.5]$, so $X$ can take only $0,1,2$.

$P(X=0)=P(10U(1-U)<1)=P(10U^2-10U+1>0)$

$=P\left(0<U<\frac12-\frac{\sqrt{60}}{20}\right)+P\left(\frac12+\frac{\sqrt{60}}{20}<U<1\right)$

$=\frac12-\frac{\sqrt{60}}{20}+1-\left(\frac12+\frac{\sqrt{60}}{20}\right)=1-\sqrt{\frac35}=0.2254$.

$P(X=1)=P(1\leq10U(1-U)<2)$

$=P(10U(1-U)<2)-P(10U(1-U)<1)$

$=P(10U^2-10U+2>0)-P(10U^2-10U+1>0)$

$=\left(\frac12-\frac{\sqrt{20}}{20}+1-\frac12-\frac{\sqrt{20}}{20}\right)-\left(1-\sqrt{\frac35}\right)$

$=\sqrt{\frac35}-\frac1{\sqrt5}=0.3274$.

$P(X=2)=P(2\leq10U(1-U))=P(10U^2-10U+2\leq0)$

$=P\left(\frac12-\frac{\sqrt{20}}{20}\leq U\leq\frac12+\frac{\sqrt{20}}{20}\right)=\frac{\sqrt{20}}{10}=\frac1{\sqrt5}=0.4472$.

\subsubsection*{Problem 4}

\term{Question.} Suppose $Y=(X-3)^2$, where $X\sim\operatorname{Binomial}(5,0.5)$. Find the pmf of $Y$.

\term{Solution.} The sample space of $X$ is $\{0,1,2,3,4,5\}$ and that of $Y$ is $\{0,1,4,9\}$. Therefore

$P(Y=0)=P(X=3)=\binom53(0.5)^5=\frac{10}{32}$,

$P(Y=1)=P(X=2\text{ or }4)=\left[\binom52+\binom54\right](0.5)^5=\frac{15}{32}$,

$P(Y=4)=P(X=1\text{ or }5)=\left[\binom51+\binom55\right](0.5)^5=\frac6{32}$,

$P(Y=9)=P(X=0)=\binom50(0.5)^5=\frac1{32}$.

\subsubsection*{Problem 5}

\term{Question.} Let $X$ have cdf

$F(x)=\begin{cases}0,&x<0,\\0.1x,&0\leq x<1,\\0.3,&1\leq x<3,\\0.2x,&3\leq x\leq5,\\1,&x>5.\end{cases}$

Determine the simulated value $x$ resulting from a uniform number $u$.

\term{Solution.}

$\begin{array}{c|c}u&\text{simulated }x\\\hline0\leq u<0.1&u=0.1x\Rightarrow x=u/0.1\\0.1\leq u\leq0.3&x=1\\0.3<u\leq0.6&x=3\\0.6<u\leq1&u=0.2x\Rightarrow x=u/0.2\end{array}$

\subsubsection*{Problem 6}

\term{Question.} Generate samples from

$f(x)=\begin{cases}\frac12x^2e^{-x},&x\geq0,\\0,&\text{otherwise},\end{cases}$

by acceptance-rejection, using

$g(x)=\begin{cases}\lambda e^{-\lambda x},&x\geq0,\\0,&\text{otherwise},\end{cases}\qquad \lambda\in(0,1)$.

Determine the best $\lambda$ minimizing the expected number of trial samples needed to generate a sample from $f$.

\term{Solution.} Let
$h(x)=\frac{f(x)}{g(x)}=\frac{x^2}{2\lambda}e^{-(1-\lambda)x}$.

Then
$h'(x)=\frac1{2\lambda}e^{-(1-\lambda)x}[2x-(1-\lambda)x^2]$.

Setting $h'(x)=0$ gives $x[2-(1-\lambda)x]=0$, so the positive stationary point is $x=2/(1-\lambda)$. Since $h'(x)>0$ before and $h'(x)<0$ after this point,

$\frac{f(x)}{g(x)}\leq h\left(\frac2{1-\lambda}\right)=\frac{2e^{-2}}{\lambda(1-\lambda)^2}$.

Thus the expected number of trial samples is $2e^{-2}/[\lambda(1-\lambda)^2]$. Define

$k(\lambda)=\frac{2e^{-2}}{\lambda(1-\lambda)^2}$.

Then
$k'(\lambda)=-\frac{2e^{-2}(1-3\lambda)}{\lambda^2(1-\lambda)^3}$.

Setting $k'(\lambda)=0$ gives $\lambda=1/3$. By the first derivative test, this is the minimum, and
$k(1/3)=\frac{2e^{-2}}{(1/3)(1-1/3)^2}=1.827$.

\subsubsection*{Problem 7}

\term{Question.} Suppose $X$ has the standard double exponential density $g(x)=\frac12e^{-|x|}$, $-\infty<x<\infty$. Show that $|X|$ has the exponential distribution with mean $1$.

\term{Solution.} Clearly $P(|X|\leq x)=0$ for $x\leq0$. For $x>0$,

$P(|X|\leq x)=P(-x\leq X\leq x)=P(X\leq x)-P(X\leq-x)$

$=\frac12\int_{-\infty}^0e^sds+\frac12\int_0^xe^{-s}ds-\frac12\int_{-\infty}^{-x}e^sds$
$=1-e^{-x}$.

Hence
$f_{|X|}(x)=\frac{d}{dx}P(|X|\leq x)=\begin{cases}0,&x\leq0,\\e^{-x},&x>0,\end{cases}$

which is exponential with mean $1$.

\subsubsection*{Problem 8}

\term{Question.} Let $\Theta$ be uniform on $[0,2\pi]$ and $S$ be exponential with mean $2$. Assume independence. Show that
$X=\sqrt S\cos\Theta,\qquad Y=\sqrt S\sin\Theta$

are independent standard normal random variables. Hint: use the two-dimensional change-of-variables formula and its Jacobian.

\term{Solution.} Let $g(x,y)$ be the joint density of $X,Y$, and $f(s,\theta)$ that of $S,\Theta$. Independence gives
$f(s,\theta)=\left(\frac1{2\pi}\right)\left(\frac12e^{-s/2}\right)=\frac1{4\pi}e^{-s/2}$.

Since $x=\sqrt s\cos\theta$ and $y=\sqrt s\sin\theta$,
$J=\left|\begin{matrix}\partial x/\partial s&\partial x/\partial\theta\\\partial y/\partial s&\partial y/\partial\theta\end{matrix}\right|$

$=\left|\begin{matrix}\frac1{2\sqrt s}\cos\theta&-\sqrt s\sin\theta\\\frac1{2\sqrt s}\sin\theta&\sqrt s\cos\theta\end{matrix}\right|=\frac12(\cos^2\theta+\sin^2\theta)=\frac12$.

Thus $f(s,\theta)=g(x,y)|J|$ and

$g(x,y)=2f(s,\theta)=\frac1{2\pi}e^{-s/2}=\frac1{2\pi}e^{-(x^2+y^2)/2}$

$=\left(\frac1{\sqrt{2\pi}}e^{-x^2/2}\right)\left(\frac1{\sqrt{2\pi}}e^{-y^2/2}\right)$.

Therefore $X$ and $Y$ are independent standard normal random variables.

\subsubsection*{Problem 9}

\term{Question.} Given

$\Sigma=\begin{pmatrix}3.21&1.82&2.18\\1.82&5.89&3.23\\2.18&3.23&8.12\end{pmatrix}$,
find its Cholesky factor $A$.

\term{Solution.} Let

$A=\begin{pmatrix}A_{11}&0&0\\A_{21}&A_{22}&0\\A_{31}&A_{32}&A_{33}\end{pmatrix}$,
and impose $AA^\top=\Sigma$. Then

$A_{11}^2=3.21\Rightarrow A_{11}=1.7916$,

$A_{21}A_{11}=1.82\Rightarrow A_{21}=1.82/1.7916=1.0159$,

$A_{31}A_{11}=2.18\Rightarrow A_{31}=2.18/1.7916=1.2168$,

$A_{21}^2+A_{22}^2=5.89\Rightarrow A_{22}=\sqrt{5.89-(1.0159)^2}=2.2041$,

$A_{31}A_{21}+A_{32}A_{22}=3.23\Rightarrow A_{32}=0.9046$,

$A_{31}^2+A_{32}^2+A_{33}^2=8.12\Rightarrow A_{33}=2.4127$.

Hence
$A=\begin{pmatrix}1.7916&0&0\\1.0159&2.2041&0\\1.2168&0.9046&2.4127\end{pmatrix}$.

\subsubsection*{Problem 10}

\term{Question.} Let $X,Y$ be arbitrary random variables and define

$\operatorname{var}(X\mid Y)=\mathbb{E}[X^2\mid Y]-\mathbb{E}[X\mid Y]^2$.

Show (a) $\operatorname{var}(X\mid Y)\geq0$ and (b)
$\operatorname{var}(X)=\mathbb{E}[\operatorname{var}(X\mid Y)]+\operatorname{var}(\mathbb{E}[X\mid Y])$.

Hints: (i) $\mathbb{E}[\mathbb{E}[f(X)\mid Y]]=\mathbb{E}[f(X)]$; (ii)

$\mathbb{E}[g(\mathbb{E}[h(X)\mid Y])k(X)\mid Y]=g(\mathbb{E}[h(X)\mid Y])\mathbb{E}[k(X)\mid Y]$.

\term{Solution.} (a) From hint (ii),

$\mathbb{E}[X\mid Y]^2=\mathbb{E}[X\mid Y]\mathbb{E}[X\mid Y]=\mathbb{E}[\mathbb{E}[X\mid Y]X\mid Y]$,

and
$\mathbb{E}[\mathbb{E}[X\mid Y]^2\mid Y]=\mathbb{E}[X\mid Y]^2$.

Therefore
$\operatorname{var}(X\mid Y)=\mathbb{E}[X^2\mid Y]-2\mathbb{E}[X\mid Y]^2+\mathbb{E}[X\mid Y]^2$

$=\mathbb{E}[X^2\mid Y]-2\mathbb{E}[\mathbb{E}[X\mid Y]X\mid Y]+\mathbb{E}[\mathbb{E}[X\mid Y]^2\mid Y]$

$=\mathbb{E}[(X-\mathbb{E}[X\mid Y])^2\mid Y]\geq0$.

(b)
$\mathbb{E}[\operatorname{var}(X\mid Y)]=\mathbb{E}[\mathbb{E}[X^2\mid Y]]-\mathbb{E}[\mathbb{E}[X\mid Y]^2]$
$=\mathbb{E}[X^2]-\mathbb{E}[\mathbb{E}[X\mid Y]^2]$,

while
$\operatorname{var}(\mathbb{E}[X\mid Y])=\mathbb{E}[\mathbb{E}[X\mid Y]^2]-\mathbb{E}[\mathbb{E}[X\mid Y]]^2$
$=\mathbb{E}[\mathbb{E}[X\mid Y]^2]-\mathbb{E}[X]^2$.

Adding gives
$\mathbb{E}[\operatorname{var}(X\mid Y)]+\operatorname{var}(\mathbb{E}[X\mid Y])=\mathbb{E}[X^2]-\mathbb{E}[X]^2=\operatorname{var}(X)$.

\subsection{Assignment 2}

\subsubsection*{Problem 1}

\term{Question.} Let $S(t)$ be the price of a non-dividend-paying asset under the risk-neutral measure $\mathbb Q$:
$\frac{dS(t)}{S(t)}=r\,dt+\sigma\,dW(t)$,

where $r$ is the continuously compounded risk-free interest rate and $W(t)$ is a standard Brownian motion under $\mathbb Q$. Derive the price of a $T$-year Asian option with payoff
$\max(S_G-K,0),\qquad S_G=\left(\prod_{i=1}^nS(t_i)\right)^{1/n}$,

over monitoring dates $0<t_1<\cdots<t_n<T$. Note: if $Y\sim LN(\mu,\beta^2)$, then

$\mathbb{E}[\max(Y-C,0)]=\mathbb{E}[Y]N(d_1)-CN(d_2)$,

$d_1=\frac{\log(\mathbb{E}[Y]/C)+\beta^2/2}{\beta},\qquad d_2=d_1-\beta$.

\term{Solution.} The SDE has solution
$S(t_i)=S(0)\exp\left[\left(r-\frac12\sigma^2\right)t_i+\sigma W(t_i)\right]$.

Therefore

\noindent\resizebox{\columnwidth}{!}{$\displaystyle\left(\prod_{i=1}^nS(t_i)\right)^{1/n}=S(0)\exp\left[\left(r-\frac12\sigma^2\right)\frac{\sum_{i=1}^nt_i}{n}+\frac\sigma n\sum_{i=1}^nW(t_i)\right]$.}

Let
$X=(W(t_1),W(t_2),\ldots,W(t_n))^\top$.

Then $X\sim N(0,\Sigma)$, where
$\Sigma=\begin{pmatrix}t_1&t_1&\cdots&t_1\\t_1&t_2&\cdots&t_2\\\vdots&\vdots&\ddots&\vdots\\t_1&t_2&\cdots&t_n\end{pmatrix}$.

Let $A=(1,1,\ldots,1)$. Since $\sum_iW(t_i)=AX$, the linear transformation property of the multivariate normal gives
$\sum_{i=1}^nW(t_i)\sim N(0,A\Sigma A^\top)$,

where
$A\Sigma A^\top=\sum_{i=1}^n(2n-2i+1)t_i$.

Thus
$\left(\prod_{i=1}^nS(t_i)\right)^{1/n}=\exp\left[\log S(0)+\left(r-\frac12\sigma^2\right)\frac{\sum_{i=1}^nt_i}{n}+\alpha Z\right]$,

where $Z\sim N(0,1)$ and
$\alpha^2=\frac{\sigma^2}{n^2}\sum_{i=1}^n(2n-2i+1)t_i$.

Hence
$S_G\sim LN\left(\log S(0)+\left(r-\frac12\sigma^2\right)\frac{\sum_{i=1}^nt_i}{n},\alpha^2\right)$,

and
$\mathbb{E}^{\mathbb Q}[S_G]=S(0)\exp\left[\left(r-\frac12\sigma^2\right)\frac{\sum_{i=1}^nt_i}{n}+\frac12\alpha^2\right]$.

Using the given lognormal result, the option price is

$V=e^{-rT}\mathbb{E}^{\mathbb Q}[\max(S_G-K,0)]$
$=e^{-rT}[\mathbb{E}^{\mathbb Q}(S_G)N(d_1)-KN(d_2)]$

$=S(0)e^{-rT}\exp\left[\left(r-\frac12\sigma^2\right)\frac{\sum_{i=1}^nt_i}{n}+\frac12\alpha^2\right]N(d_1)-Ke^{-rT}N(d_2)$,

where
$d_1=\frac{\log(\mathbb{E}^{\mathbb Q}[S_G]/K)+\alpha^2/2}{\alpha}$

$=\frac{\log(S(0)/K)+\left(r-\frac12\sigma^2\right)\frac{\sum_{i=1}^nt_i}{n}+\alpha^2}{\alpha},\qquad d_2=d_1-\alpha$.

\subsubsection*{Problem 2}

\term{Question.} The dynamics of $S_1,S_2$ are
$\frac{dS_i(t)}{S_i(t)}=\mu_i\,dt+\sigma_i\,dX_i(t),\qquad i=1,2$,

where $X_i$ are standard Brownian motions and $dX_1(t)dX_2(t)=\rho\,dt$. Suppose $S_1(0)=5$, $S_2(0)=6$, $\mu_1=0.03$, $\mu_2=0.05$, $\sigma_1=0.3$, $\sigma_2=0.2$, and $\rho=-0.2$. Find $\operatorname{Cov}(S_1(0.5),S_2(0.5))$.

\term{Solution.} The exact solutions are

$S_i(t)=S_i(0)\exp\left[\left(\mu_i-\frac12\sigma_i^2\right)t+\sigma_iX_i(t)\right],\qquad i=1,2$.

Therefore

$\operatorname{Cov}(S_1(t),S_2(t))=\mathbb{E}[S_1(t)S_2(t)]-\mathbb{E}[S_1(t)]\mathbb{E}[S_2(t)]$

$=S_1(0)S_2(0)\exp\left[\left(\mu_1+\mu_2-\frac12\sigma_1^2-\frac12\sigma_2^2\right)t\right]$

$\qquad\times\mathbb{E}\left[e^{\sigma_1X_1(t)+\sigma_2X_2(t)}\right]-S_1(0)S_2(0)e^{(\mu_1+\mu_2)t}$.

At $t=0.5$ this becomes

$30e^{0.0075}\mathbb{E}[e^{0.3X_1(0.5)+0.2X_2(0.5)}]-30e^{0.04}$.

Let $Y=(Y_1,Y_2)^\top\sim N(0,\Sigma)$ with
$\Sigma=\begin{pmatrix}1&-0.2\\-0.2&1\end{pmatrix}$.

Then
$\mathbb{E}[e^{0.3X_1(0.5)+0.2X_2(0.5)}]=\mathbb{E}[e^{0.3\sqrt{0.5}Y_1+0.2\sqrt{0.5}Y_2}]$

$=\exp\left[\frac12\{0.3^2(0.5)+0.2^2(0.5)-2(0.2)(0.3)(0.2)(0.5)\}\right]=e^{0.0265}$.

Hence
$\operatorname{Cov}(S_1(0.5),S_2(0.5))=30e^{0.0075}e^{0.0265}-30e^{0.04}=-0.1868$.

\subsubsection*{Problem 3}

\term{Question.} Show that the optimal value of $b$ in the control variate method is $\operatorname{Cov}(X,Y)/\operatorname{var}(X)$.

\term{Solution.} The control-variate estimator is

$\overline{Y}(b)=\frac1n\sum_{i=1}^n[Y_i-b(X_i-\mathbb{E}[X])]$.

Its variance is
$\operatorname{var}(\overline{Y}(b))=\frac{\sigma_Y^2-2b\rho_{XY}\sigma_X\sigma_Y+b^2\sigma_X^2}{n}=:\sigma^2(b)$.

The optimal value is
$b^*=\arg\min_b\sigma^2(b)$.

Differentiating,
$\frac{d}{db}\sigma^2(b)=\frac{-2\rho_{XY}\sigma_X\sigma_Y+2b\sigma_X^2}{n}=0$,

so
$b^*=\frac{\rho_{XY}\sigma_X\sigma_Y}{\sigma_X^2}=\frac{\operatorname{Cov}(X,Y)}{\operatorname{var}(X)}$.

\subsubsection*{Problem 4}

\term{Question.} Suppose $X$ is Bernoulli with parameter $p$. Show that if $2p>1$, then an antithetic variable of $X$ is
$Y=1-XZ$,

where $Z$ is Bernoulli with parameter $(1-p)/p$, independent of $X$. Show both that $X,Y$ have the same distribution and that $\operatorname{Cov}(X,Y)<0$.

\term{Solution.}

$P(Y=0)=P(X=1\text{ and }Z=1)=p\frac{1-p}{p}=1-p$,
and

$P(Y=1)=P(X=1,Z=0)+P(X=0,Z=1)+P(X=0,Z=0)$

$=p\left(1-\frac{1-p}{p}\right)+(1-p)\frac{1-p}{p}+(1-p)\left(1-\frac{1-p}{p}\right)=p$.

Thus $X$ and $Y$ have the same distribution. Also,

$\operatorname{Cov}(X,Y)=\mathbb{E}[XY]-\mathbb{E}[X]\mathbb{E}[Y]$
$=\mathbb{E}[X(1-XZ)]-\mathbb{E}[X]\mathbb{E}[1-XZ]$

$=\mathbb{E}[X]-\mathbb{E}[X^2Z]-\mathbb{E}[X]+\mathbb{E}[X]\mathbb{E}[XZ]$
$=-\mathbb{E}[XZ]+\mathbb{E}[X]\mathbb{E}[XZ]$

$=-(1-\mathbb{E}[X])\mathbb{E}[X]\mathbb{E}[Z]$
$=-\left(1-p\right)p\left(\frac{1-p}{p}\right)=-(1-p)^2<0$.

Hence $Y$ is antithetic to $X$.

\subsubsection*{Problem 5}

\term{Question.} Use antithetic variates to estimate $P(Z\geq0.8)$ for standard normal $Z$. (a) Write the estimator in terms of i.i.d. $Z_1,\ldots,Z_n\sim N(0,1)$. (b) Compare it with plain Monte Carlo using $2n$ independent replications and determine the percentage variance reduction.

\term{Solution.} (a) Let $Y_i=\mathbf1_{\{Z_i\geq0.8\}}$. Since $f(x)=\mathbf1_{\{x\geq0.8\}}$ is monotone, the antithetic variable is $\widetilde Y_i=\mathbf1_{\{-Z_i\geq0.8\}}=\mathbf1_{\{Z_i\leq-0.8\}}$. Hence

$\widehat Y_{AV}=\frac1{2n}\sum_{i=1}^n(Y_i+\widetilde Y_i)$.

(b)
$\operatorname{var}(\widehat Y_{AV})=\frac1{4n^2}\sum_{i=1}^n\operatorname{var}(Y_i+\widetilde Y_i)$
$=\frac1{4n}\{2\operatorname{var}(Y_1)+2\operatorname{cov}(Y_1,\widetilde Y_1)\}$

$=\frac1{2n}\{\operatorname{var}(Y_1)+\operatorname{cov}(Y_1,\widetilde Y_1)\}$.

For plain Monte Carlo with $2n$ draws,
$\operatorname{var}(\widehat Y)=\frac1{4n^2}\sum_{i=1}^{2n}\operatorname{var}(Y_i)=\frac1{2n}\operatorname{var}(Y_1)$

$=\frac1{2n}\{\mathbb{E}[Y_1^2]-\mathbb{E}[Y_1]^2\}$
$=\frac1{2n}P(Z\geq0.8)[1-P(Z\geq0.8)]$

$=\frac1{2n}(1-0.7881)(0.7881)=\frac{0.0835}{n}$.

The magnitude of variance change is
$\operatorname{var}(\widehat Y_{AV})-\operatorname{var}(\widehat Y)=\frac1{2n}\operatorname{cov}(Y_1,\widetilde Y_1)$.

Now
$\operatorname{cov}(Y_1,\widetilde Y_1)=\mathbb{E}[\mathbf1_{\{Z\geq0.8\}}\mathbf1_{\{Z\leq-0.8\}}]-\mathbb{E}[\mathbf1_{\{Z\geq0.8\}}]\mathbb{E}[\mathbf1_{\{Z\leq-0.8\}}]$

$=0-P(Z\geq0.8)P(Z\leq-0.8)=-(0.2119)^2=-0.0449$.

Therefore
$\operatorname{var}(\widehat Y_{AV})-\operatorname{var}(\widehat Y)=\frac{-0.0449}{2n}=-\frac{0.0225}{n}$.

The percentage variance change is
$\frac{-0.0225/n}{0.0835/n}\times100\%=-26.9461\%$.

\subsubsection*{Problem 6}

\term{Question.} Let $U\sim\operatorname{Unif}(0,1)$ and estimate $\mathbb{E}[e^U]$ by stratified sampling with two equal strata. Let $V_{11},\ldots,V_{1n}$ be independent uniform variables on $[0,0.5]$, and $V_{21},\ldots,V_{2n}$ independent uniform variables on $[0.5,1]$. (a) Write the proportional-allocation estimator. (b) Evaluate the magnitude of variance reduction. (c) Find the optimal allocation and corresponding estimator.

\term{Solution.} Let $Y=e^U$, $A_1=[0,1/2]$, $A_2=(1/2,1]$. The stratum probabilities are $p_1=p_2=1/2$ and proportional allocation gives $n_1=n_2=n$.

(a) The estimator is
$\widehat Y=\frac1{2n}\sum_{i=1}^2\sum_{j=1}^nY_{ij}$,
where $Y_{1j}=e^{V_{1j}}$ and $Y_{2j}=e^{V_{2j}}$.

(b) The conditional means are

$\mu_1=\mathbb{E}[Y_{1j}]=\mathbb{E}[e^{V_{1j}}]=2\int_0^{0.5}e^x\,dx=2(\sqrt e-1)$,

$\mu_2=\mathbb{E}[Y_{2j}]=\mathbb{E}[e^{V_{2j}}]=2\int_{0.5}^1e^x\,dx=2(e-\sqrt e)$.

The magnitude of variance reduction is

$\frac1n\left[\sum_{i=1}^2p_i\mu_i^2-\left(\sum_{i=1}^2p_i\mu_i\right)^2\right]$

$=\frac1n\left[\frac12\{2(\sqrt e-1)\}^2+\frac12\{2(e-\sqrt e)\}^2-(e-1)^2\right]$
$=\frac{0.0886}{n}$.

(c) The optimal allocation is
$q_i^*=\frac{p_i\sigma_i}{\sum_{l=1}^2p_l\sigma_l},\qquad \sigma_i^2=\operatorname{var}(Y_{ij})$.

For stratum 1,
$\sigma_1^2=\operatorname{var}(e^{V_{1j}})=\mathbb{E}[e^{2V_{1j}}]-\mathbb{E}[e^{V_{1j}}]^2$

$=2\int_0^{0.5}e^{2x}\,dx-4(\sqrt e-1)^2=e-1-4(\sqrt e-1)^2=0.0349$.

For stratum 2,
$\sigma_2^2=2\int_{0.5}^1e^{2x}\,dx-4(e-\sqrt e)^2=e^2-e-4(e-\sqrt e)^2=0.0949$.

Thus
$q_1^*=\frac{0.5\sqrt{0.0349}}{0.5\sqrt{0.0349}+0.5\sqrt{0.0949}}=0.3775$,

$q_2^*=\frac{0.5\sqrt{0.0949}}{0.5\sqrt{0.0349}+0.5\sqrt{0.0949}}=0.6225$.

The corresponding estimator is
$\widehat Y^*=\frac1{4n}\left(\frac1{q_1^*}\sum_{j=1}^{2nq_1^*}Y_{1j}+\frac1{q_2^*}\sum_{j=1}^{2nq_2^*}Y_{2j}\right)$.

\subsubsection*{Problem 7}

\term{Question.} Let $X\sim\operatorname{Binomial}(m,p)$ with
$P(X=k)=\frac{m!}{(m-k)!k!}p^k(1-p)^{m-k}$.

Write the importance-sampling estimator for $P(X\geq m\alpha)$, where $\alpha\in(0,1)$, using the alternative distribution $\operatorname{Binomial}(m,\widetilde p)$.

\term{Solution.} Rewrite
$P(X\geq m\alpha)=\mathbb{E}[\mathbf1_{\{X\geq m\alpha\}}]$.

Let $Y_i\sim\operatorname{Binomial}(m,\widetilde p)$ independently. Then

$\widehat p_{IS}=\frac1n\sum_{i=1}^n\mathbf1_{\{Y_i\geq m\alpha\}}\frac{\binom m{Y_i}p^{Y_i}(1-p)^{m-Y_i}}{\binom m{Y_i}\widetilde p^{Y_i}(1-\widetilde p)^{m-Y_i}}$

$=\frac1n\sum_{i=1}^n\mathbf1_{\{Y_i\geq m\alpha\}}\left(\frac p{\widetilde p}\right)^{Y_i}\left(\frac{1-p}{1-\widetilde p}\right)^{m-Y_i}$.

\subsubsection*{Problem 8}

\term{Question.} Let $X\sim N(0,1)$ and estimate $\mathbb{E}[\mathbf1_{\{X\geq b\}}]$, $b>0$, by importance sampling with $g$ the density of $N(\theta,1)$. (a) Let $\theta=x^*$, where $x^*$ maximizes $\mathbf1_{\{x\geq b\}}f(x)$ for the $N(0,1)$ density $f$. Determine $\theta$. (b) Using $g=N(x^*,1)$, write the estimator in terms of $X_i\sim N(x^*,1)$.

\term{Solution.} (a)
$\mathbf1_{\{x\geq b\}}f(x)=\begin{cases}0,&x<b,\\f(x),&x\geq b.\end{cases}$

Since $f$ is nonnegative and decreasing on $[0,\infty)$, $x^*=b$.

(b) Let $g$ be the density of $N(b,1)$:
$g(x)=\frac1{\sqrt{2\pi}}e^{-(x-b)^2/2}$.

Then
$\widehat p_{IS}=\frac1n\sum_{i=1}^n\mathbf1_{\{X_i\geq b\}}\frac{f(X_i)}{g(X_i)}$
$=\frac1n\sum_{i=1}^n\mathbf1_{\{X_i\geq b\}}\exp\left[-\frac{X_i^2}{2}+\frac{(X_i-b)^2}{2}\right]$

$=\frac1n\sum_{i=1}^n\mathbf1_{\{X_i\geq b\}}\exp\left(-bX_i+\frac12b^2\right)$.

\subsubsection*{Problem 9}

\term{Question.} Determine the exponential tilt family when (a) $X$ is Bernoulli with parameter $p$ and (b) $X$ is exponential with mean $1/\lambda$.

\term{Solution.} (a) The Bernoulli pmf is
$p(x)=\begin{cases}1-p,&x=0,\\p,&x=1,\\0,&\text{otherwise}.\end{cases}$

The exponential tilt is
$p_\theta(x)=\frac{e^{\theta x}p(x)}{\mathbb{E}[e^{\theta X}]}=\frac{e^{\theta x}p(x)}{pe^\theta+1-p}$,

so

$p_\theta(x)=\begin{cases}\dfrac{1-p}{pe^\theta+1-p},&x=0,\\\dfrac{pe^\theta}{pe^\theta+1-p},&x=1,\\0,&\text{otherwise}.\end{cases}$

(b) For $f(x)=\lambda e^{-\lambda x}\mathbf1_{\{x\geq0\}}$,

$f_\theta(x)=\frac{e^{\theta x}f(x)}{\mathbb{E}[e^{\theta X}]}=\frac{e^{\theta x}\lambda e^{-\lambda x}}{\lambda/(\lambda-\theta)}$
$=(\lambda-\theta)e^{-(\lambda-\theta)x},\qquad x\geq0,\quad\theta<\lambda$.

\subsubsection*{Problem 10}

\term{Question.} Let $X$ have density $f$, let $\{f_\theta\}$ be an exponential tilt family, and let $H(\theta)=\log\mathbb{E}[e^{\theta X}]$. Show that the general iterative cross-entropy update is

$H'(\widehat\theta_{j+1})=\frac{\sum_{k=1}^Nh(Y_k)\ell(Y_k;\widehat\theta_j)Y_k}{\sum_{k=1}^Nh(Y_k)\ell(Y_k;\widehat\theta_j)}$,

where $Y_1,\ldots,Y_N$ are i.i.d. pilot samples from $f_{\widehat\theta_j}$.

\term{Solution.} The parameter $\widehat\theta$ solves
$0=\frac1N\sum_{k=1}^Nh(X_k)\frac{\partial}{\partial\theta}\log f_\theta(X_k)$.

Since
$f_\theta(x)=\frac{e^{\theta x}f(x)}{\mathbb{E}[e^{\theta X}]}$,
$\log f_\theta(X_k)=\log f(X_k)+\theta X_k-H(\theta)$.

Therefore
$0=\frac1N\sum_{k=1}^Nh(X_k)[X_k-H'(\theta)]$,

and hence
$H'(\widehat\theta)=\frac{\sum_{k=1}^Nh(X_k)X_k}{\sum_{k=1}^Nh(X_k)}$.

Using pilot samples from $f_{\widehat\theta_j}$ and the likelihood ratio $\ell(Y_k;\widehat\theta_j)$, the next update satisfies
$0=\frac1N\sum_{k=1}^Nh(Y_k)\ell(Y_k;\widehat\theta_j)\left.\frac{\partial}{\partial\theta}\log f_\theta(Y_k)\right|_{\theta=\widehat\theta_{j+1}}$

$=\frac1N\sum_{k=1}^Nh(Y_k)\ell(Y_k;\widehat\theta_j)[Y_k-H'(\widehat\theta_{j+1})]$.

Thus
$H'(\widehat\theta_{j+1})=\frac{\sum_{k=1}^Nh(Y_k)\ell(Y_k;\widehat\theta_j)Y_k}{\sum_{k=1}^Nh(Y_k)\ell(Y_k;\widehat\theta_j)}$.

\subsection{Assignment 3}

\subsubsection*{Problem 1}

\term{Question.} Let $X_1(t),X_2(t)$ satisfy
$dX_1(t)=a_1X_1(t)dt+\sigma_1X_1(t)dW_1(t)$,

$dX_2(t)=a_2X_2(t)dt+\sigma_2X_2(t)dW_2(t)$,
where $a_1,a_2,\sigma_1,\sigma_2$ are constants and $dW_1(t)dW_2(t)=\rho\,dt$. Let $(R_1,V_1),\ldots,(R_m,V_m)$ be i.i.d. samples from $N(0,I_2)$. On $0=t_0<t_1<\cdots<t_m$, write the Euler scheme in terms of these samples.

\term{Solution.} Let $\widehat X_1,\widehat X_2$ be the time-discretized approximations. Let $Z_1,Z_2$ be independent standard Brownian motions. Since $W_1,W_2$ are correlated,

$dW_1(t)=dZ_1(t),\qquad dW_2(t)=\rho\,dZ_1(t)+\sqrt{1-\rho^2}\,dZ_2(t)$.

Set $\widehat X_1(0)=X_1(0)$ and $\widehat X_2(0)=X_2(0)$. For $i=0,1,\ldots,m-1$,

$\widehat X_1(t_{i+1})=\widehat X_1(t_i)+a_1\widehat X_1(t_i)(t_{i+1}-t_i)+\sigma_1\widehat X_1(t_i)\sqrt{t_{i+1}-t_i}\,R_{i+1}$,

$\widehat X_2(t_{i+1})=\widehat X_2(t_i)+a_2\widehat X_2(t_i)(t_{i+1}-t_i)$

$\qquad+\sigma_2\widehat X_2(t_i)\sqrt{t_{i+1}-t_i}\left(\rho R_{i+1}+\sqrt{1-\rho^2}V_{i+1}\right)$.

\subsubsection*{Problem 2}

\term{Question.} Let $Z_1,\ldots,Z_m$ be i.i.d. standard normal variables. On $0=t_0<t_1<\cdots<t_m$, write the Euler and Milstein schemes for

$dX(t)=rX(t)dt+\sigma\sqrt{\sigma_0^2+\sigma_1^2X^2(t)}\,dW(t)$,

where all constants are strictly positive.

\term{Solution.} Let $\widehat X$ be the time-discretized approximation and set $\widehat X(0)=X(0)$.

\term{Euler scheme.} For $i=0,1,\ldots,m-1$,

$\widehat X(t_{i+1})=\widehat X(t_i)+r\widehat X(t_i)(t_{i+1}-t_i)$
$+\sqrt{\sigma_0^2+\sigma_1^2\widehat X^2(t_i)}\sqrt{t_{i+1}-t_i}\,Z_{i+1}$.

\term{Milstein scheme.} For $i=0,1,\ldots,m-1$,

$\widehat X(t_{i+1})=\widehat X(t_i)+r\widehat X(t_i)(t_{i+1}-t_i)$
$+\sqrt{\sigma_0^2+\sigma_1^2\widehat X^2(t_i)}\sqrt{t_{i+1}-t_i}\,Z_{i+1}$

$\qquad+\frac12\left(\frac{2\sigma_1^2\widehat X(t_i)}{2\sqrt{\sigma_0^2+\sigma_1^2\widehat X^2(t_i)}}\right)\sqrt{\sigma_0^2+\sigma_1^2\widehat X^2(t_i)}(t_{i+1}-t_i)(Z_{i+1}^2-1)$

$=\widehat X(t_i)+r\widehat X(t_i)(t_{i+1}-t_i)+\sqrt{\sigma_0^2+\sigma_1^2\widehat X^2(t_i)}\sqrt{t_{i+1}-t_i}\,Z_{i+1}$

$\qquad+\frac12\sigma_1^2\widehat X(t_i)(t_{i+1}-t_i)(Z_{i+1}^2-1)$.

\subsubsection*{Problem 3}

\term{Question.} Find the Lamperti transforms of: (a)

$dX(t)=a(b-X(t))dt+\sigma\sqrt{X(t)}dW(t)$;

(b)
$dX(t)=(r_0+r_1X(t))dt+\sigma X(t)dW(t)$.

All constants are strictly positive.

\term{Solution.} (a)
$F(x)=\int\frac1{\sigma\sqrt x}\,dx=\frac2\sigma\sqrt x,\qquad Y(t)=F(X(t))=\frac2\sigma\sqrt{X(t)},$

and $X(t)=\sigma^2Y^2(t)/4$. By It\^o's lemma,

$dY(t)=\frac{\partial Y}{\partial X}dX(t)+\frac12\frac{\partial^2Y}{\partial X^2}(dX(t))^2$
$=\frac2\sigma\left(\frac12\right)\frac1{\sqrt{X(t)}}[a(b-X(t))dt+\sigma\sqrt{X(t)}dW(t)]$
$+\frac12\left(-\frac1{2\sigma}\right)\frac1{\sqrt{X^3(t)}}[\sigma^2X(t)dt]$

$=a\left(\frac{b}{\sigma\sqrt{X(t)}}-\frac{\sqrt{X(t)}}\sigma-\frac\sigma{4a\sqrt{X(t)}}\right)dt+dW(t)$.

(b)
$F(x)=\int\frac1{\sigma x}\,dx=\frac1\sigma\log x,\qquad Y(t)=F(X(t))=\frac1\sigma\log X(t),$

and $X(t)=e^{\sigma Y(t)}$. By It\^o's lemma,

$dY(t)=\frac1{\sigma X(t)}[(r_0+r_1X(t))dt+\sigma X(t)dW(t)]-\frac1{2\sigma X^2(t)}[\sigma^2X^2(t)dt]$

$=\left(\frac{r_0}{\sigma X(t)}-\frac{r_1}\sigma-\frac\sigma2\right)dt+dW(t)$.

\subsubsection*{Problem 4}

\term{Question.} Assume a GBM stock under the risk-neutral measure,

$dS(t)=rS(t)dt+\sigma S(t)dW(t)$.

Using pathwise differentiation, the vega of a standard European call with maturity $T$ and strike $K$ can be written
$\mathbb{E}\left[e^{-rT}\mathbf1_{\{S(T)>K\}}f(S(0),S(T))\right].$

Find $f(S(0),S(T))$.

\term{Solution.} Let $Y=e^{-rT}[S(T)-K]^+$, where

$S(T)=S(0)e^{(r-\sigma^2/2)T+\sigma\sqrt T Z},\qquad Z\sim N(0,1)$.

Take the parameter $\theta$ to be $\sigma$. The pathwise vega is
$\nu=\mathbb{E}\left[\frac{dY}{d\sigma}\right]$.

By the chain rule,
$\frac{dY}{d\sigma}=\frac{dY}{dS(T)}\frac{dS(T)}{d\sigma}$.

The first factor is
$\frac{dY}{dS(T)}=e^{-rT}\mathbf1_{\{S(T)>K\}}$.

For the second,
$\frac{dS(T)}{d\sigma}=S(0)e^{(r-\sigma^2/2)T+\sigma\sqrt T Z}(-\sigma T+\sqrt T Z)$.

Eliminating $Z$ gives
$\frac{dS(T)}{d\sigma}=\frac{S(T)}\sigma\left[\log\frac{S(T)}{S(0)}-rT-\frac12\sigma^2T\right]$.

Hence
$\nu=\mathbb{E}\left[e^{-rT}\mathbf1_{\{S(T)>K\}}\frac{S(T)}\sigma\left(\log\frac{S(T)}{S(0)}-rT-\frac12\sigma^2T\right)\right]$,

so
$f(S(0),S(T))=\frac{S(T)}\sigma\left(\log\frac{S(T)}{S(0)}-rT-\frac12\sigma^2T\right)$.

\subsubsection*{Problem 5}

\term{Question.} Under the same risk-neutral GBM, consider an arithmetic-average option with strike $K$ and payoff
$(S_A-K)^+,\qquad S_A=\frac1m\sum_{i=1}^mS(t_i),$

where $t_i=ih$, $h=T/m$. Using the likelihood ratio method, write a time-$0$ vega estimate in terms of i.i.d. $Z_1,\ldots,Z_m\sim N(0,1)$.

\term{Solution.} 
Let $X\sim g_\theta(x)$, $V(\theta)=E_\theta[h(X)]
=\int h(x)g_\theta(x)\,dx$. 

$
\frac{\partial V}{\partial\theta}
=
E_\theta\left[
h(X)
\frac{\partial_\theta g_\theta(X)}
{g_\theta(X)}
\right]
$
$
=
E_\theta\left[
h(X)
\frac{\partial}{\partial\theta}\log g_\theta(X)
\right]
$

Under GBM,

$S(t_i)=S(t_{i-1})\exp\left[\left(r-\frac12\sigma^2\right)h+\sigma\sqrt hZ_i\right],\qquad i=1,\ldots,m$.

Let $g(x_1,\ldots,x_m)$ be the joint density of $S(t_1),\ldots,S(t_m)$. By the Markov property,
$g(x_1,\ldots,x_m)=g_1(x_1\mid x_0)g_2(x_2\mid x_1)\cdots g_m(x_m\mid x_{m-1}),\qquad x_0=S(0)$,

where
$g_j(x_j\mid x_{j-1})=\frac1{x_j\sigma\sqrt h}f(\xi_j)$,
$\xi_j=\frac{\log(x_j/x_{j-1})-(r-\sigma^2/2)h}{\sigma\sqrt h}$,

and $f$ is the standard normal density. Hence

$\frac\partial{\partial\sigma}\log g(S(t_1),\ldots,S(t_m))=\sum_{j=1}^m\frac\partial{\partial\sigma}\log g_j(S(t_j)\mid S(t_{j-1}))$.

For each term,
$\frac\partial{\partial\sigma}\log g_j(S(t_j)\mid S(t_{j-1}))$

$=\frac\partial{\partial\sigma}\left[-\log S(t_j)-\log\sigma-\frac12\log h-\frac12\log(2\pi)-\frac12\xi_j^2\right]$
$=-\frac1\sigma-\xi_j\frac{\partial\xi_j}{\partial\sigma}$.

Since
$\frac{\partial\xi_j}{\partial\sigma}=-\frac{Z_j}\sigma+\sqrt h$,
we have

$\frac\partial{\partial\sigma}\log g_j(S(t_j)\mid S(t_{j-1}))=-\frac1\sigma-Z_j\left(-\frac{Z_j}\sigma+\sqrt h\right)$
$=\frac{Z_j^2-1}\sigma-\sqrt hZ_j$.

Therefore
$\frac\partial{\partial\sigma}\log g(S(t_1),\ldots,S(t_m))=\sum_{j=1}^m\left(\frac{Z_j^2-1}\sigma-\sqrt hZ_j\right)$.

The likelihood-ratio vega estimator is

$e^{-rT}(S_A-K)^+\sum_{j=1}^m\left(\frac{Z_j^2-1}\sigma-\sqrt hZ_j\right)$.

\subsubsection*{Problem 6}

\term{Question.} A $1.5$-year American put has strike $20$, exercise dates $0.5,1,1.5$, and continuously compounded risk-free rate $6\%$. Six risk-neutral stock paths are

$\begin{array}{c|ccc}\text{Path}&t=0.5&t=1&t=1.5\\\hline1&21.74&22.82&17.16\\2&17.98&15.84&25.89\\3&28.07&38.63&37.94\\4&10.38&13.26&10.67\\5&21.98&14.68&11.05\\6&18.88&22.78&34.57\end{array}$

Using least squares and assuming continuation value at $t_i=0.5i$, $i=1,2$, is linear in $S(t_i)$, calculate the time-$0$ value.

\term{Solution.} Terminal payoffs are

$\begin{array}{c|cc}\text{Path}&S(1.5)&Y_3=(20-S(1.5))^+\\\hline1&17.16&2.84\\2&25.89&0\\3&37.94&0\\4&10.67&9.33\\5&11.05&8.95\\6&34.57&0\end{array}$

At $t=1$, model expected continuation payoff by a linear function $f_2(S(1))$ and use the in-the-money data:

$\begin{array}{c|ccc}\text{Path}&Y_3e^{-0.03}&S(1)&\text{ITM?}\\\hline1&2.76&22.82&\text{No}\\2&0&15.84&\text{Yes}\\3&0&38.63&\text{No}\\4&9.05&13.26&\text{Yes}\\5&8.69&14.68&\text{Yes}\\6&0&22.78&\text{No}\end{array}$

The least-squares regression gives

$\mathbb{E}[Y_3e^{-0.03}\mid S(1)]=f_2(S(1))=-3.3878S(1)+55.353$.

Define
$Y_2=\begin{cases}20-S(1),&20-S(1)\geq f_2(S(1)),\\e^{-0.03}Y_3,&\text{otherwise}.\end{cases}$

The exercise comparison is

$\begin{array}{c|ccc|c}\text{Path}&20-S(1)&f_2(S(1))&Y_3e^{-0.03}&Y_2\\\hline1&0&0&2.76&2.76\\2&4.16&1.69&0&4.16\\3&0&0&0&0\\4&6.74&10.43&9.05&9.05\\5&5.32&5.62&8.69&8.69\\6&0&0&0&0\end{array}$

At $t=0.5$, use

$\begin{array}{c|ccc}\text{Path}&Y_2e^{-0.03}&S(0.5)&\text{ITM?}\\\hline1&2.68&21.74&\text{No}\\2&4.04&17.98&\text{Yes}\\3&0&28.07&\text{No}\\4&8.78&10.38&\text{Yes}\\5&8.43&21.98&\text{No}\\6&0&18.88&\text{Yes}\end{array}$

to obtain
$\mathbb{E}[Y_2e^{-0.03}\mid S(0.5)]=f_1(S(0.5))=-0.8736S(0.5)+18.0303$.

Define
$Y_1=\begin{cases}20-S(0.5),&20-S(0.5)\geq f_1(S(0.5)),\\e^{-0.03}Y_2,&\text{otherwise}.\end{cases}$

Then

$\begin{array}{c|ccc|c}\text{Path}&20-S(0.5)&f_1(S(0.5))&Y_2e^{-0.03}&Y_1\\\hline1&0&0&2.68&2.68\\2&2.02&2.32&4.04&4.04\\3&0&0&0&0\\4&9.62&8.96&8.78&9.62\\5&0&0&8.43&8.43\\6&1.12&1.54&0&0\end{array}$

Finally, the time-$0$ estimate is the average of $e^{-0.03}Y_1$:

$\widehat V_0=\frac16\sum_{i=1}^6e^{-0.03}Y_{1i}=4.0063$.

\subsubsection*{Problem 7}

\term{Question.} Consider a portfolio: long $2$ calls on $S_1$, long $5$ units of $S_1$, short $3$ puts on $S_2$, and short $4$ units of $S_2$. The call has delta $0.56$ and gamma $0.08$; the put has delta $-0.27$ and gamma $0.07$. Also

$\Delta S=(\Delta S_1,\Delta S_2)^\top\sim N(\mu,\Sigma),\qquad \mu=\binom00,\quad\Sigma=\begin{pmatrix}0.5&0.18\\0.18&0.3\end{pmatrix}$.

Let $V(t,S_1,S_2)$ be the portfolio value and neglect $\Delta t$. (a) Write the delta-gamma approximation $Q$ to the loss $L=-\Delta V$ in terms of $\Delta S$. (b) Diagonalize $Q$ and express it using independent standard normals $Z_1,Z_2$.

\term{Solution.} (a) The portfolio value is

$V(t,S_1,S_2)=2C(S_1)-3P(S_2)-4S_2+5S_1$.

The delta-gamma approximation is $\Delta V$

$\approx\frac{\partial V}{\partial S_1}\Delta S_1+\frac{\partial V}{\partial S_2}\Delta S_2+\frac12\frac{\partial^2V}{\partial S_1^2}(\Delta S_1)^2+\frac{\partial^2V}{\partial S_1\partial S_2}\Delta S_1\Delta S_2+\frac12\frac{\partial^2V}{\partial S_2^2}(\Delta S_2)^2$.

Therefore
$Q=-\Delta V=-\left(2\frac{\partial C}{\partial S_1}+5\right)\Delta S_1-\left(-3\frac{\partial P}{\partial S_2}-4\right)\Delta S_2$

$\qquad-\frac12\left(2\frac{\partial^2C}{\partial S_1^2}\right)(\Delta S_1)^2-\frac12\left(-3\frac{\partial^2P}{\partial S_2^2}\right)(\Delta S_2)^2$

$=-6.12\Delta S_1+3.19\Delta S_2-0.08(\Delta S_1)^2+0.105(\Delta S_2)^2$

$=\alpha^\top\Delta S+\Delta S^\top\Gamma\Delta S$,

where
$\alpha^\top=(-6.12,3.19),\qquad \Gamma=\begin{pmatrix}-0.08&0\\0&0.105\end{pmatrix}$.

(b) Let the Cholesky factor of $\Sigma$ be
$\widetilde C=\begin{pmatrix}\widetilde C_{11}&0\\\widetilde C_{12}&\widetilde C_{22}\end{pmatrix}$.

From $\widetilde C\widetilde C^\top=\Sigma$,
$\widetilde C_{11}=\sqrt{0.5}=0.7071,\qquad \widetilde C_{12}=0.18/0.7071=0.2546,$

$\widetilde C_{22}=\sqrt{0.3-(0.2546)^2}=0.485$.

Then
$A=\widetilde C^\top\Gamma\widetilde C=\begin{pmatrix}-0.0332&0.013\\0.013&0.0247\end{pmatrix}$.

The eigenvalues are $\lambda_1=-0.0360$ and $\lambda_2=0.0275$. Corresponding normalized eigenvectors are
$v_1=\binom{0.9778}{-0.2095},\qquad v_2=\binom{0.2095}{0.9778}$.

Let
$\Lambda=\begin{pmatrix}-0.036&0\\0&0.0424\end{pmatrix},\, U=(v_1\ v_2)=\begin{pmatrix}0.9778&0.2095\\-0.2095&0.9778\end{pmatrix}$.

Set
$C=\widetilde CU=\begin{pmatrix}0.6914&0.1481\\0.1473&0.5276\end{pmatrix}$,

and $\Delta S=CZ$. Then
$Q=\alpha^\top CZ+Z^\top(C^\top\Gamma C)Z$

$=-3.7614Z_1+0.7764Z_2-0.036Z_1^2+0.0275Z_2^2$.

\subsubsection*{Problem 8}

\term{Question.} Define $P_\theta$ by the likelihood ratio in the lecture notes. Let $Z$ be standard multivariate normal under $P$. Show that under $P_\theta$,
$Z\sim N(\mu(\theta),\Sigma(\theta))$,

where $\Sigma(\theta)$ is diagonal with
$\mu_j(\theta)=\frac{\theta b_j}{1-2\lambda_j\theta},\qquad \sigma_j^2(\theta)=\frac1{1-2\lambda_j\theta}$.

\term{Solution.} Consider the mgf of $Z$ under $P_\theta$:
$\mathbb{E}_\theta[e^{\beta^\top Z}]=\mathbb{E}[e^{\theta Q-\psi(\theta)+\beta^\top Z}]$.

For $Q=a+\sum_{j=1}^m(b_jZ_j+\lambda_jZ_j^2)$ and

$\psi(\theta)=a\theta+\frac12\sum_{j=1}^m\left(\frac{\theta^2b_j^2}{1-2\theta\lambda_j}-\log(1-2\theta\lambda_j)\right)$,

we obtain
$\mathbb{E}_\theta[e^{\beta^\top Z}]$

$=\exp\left[-\frac12\sum_{j=1}^m\left(\frac{\theta^2b_j^2}{1-2\theta\lambda_j}-\log(1-2\theta\lambda_j)\right)\right]$

$\qquad\times\mathbb{E}\left[\exp\left(\theta\sum_{j=1}^m(b_jZ_j+\lambda_jZ_j^2)+\sum_{j=1}^m\beta_jZ_j\right)\right]$

$=\exp\left[-\frac12\sum_{j=1}^m\left(\frac{\theta^2b_j^2}{1-2\theta\lambda_j}-\log(1-2\theta\lambda_j)\right)\right]$

$\qquad\times\exp\left[\frac12\sum_{j=1}^m\left(\frac{(\theta b_j+\beta_j)^2}{1-2\theta\lambda_j}-\log(1-2\theta\lambda_j)\right)\right]$

$=\exp\left[\sum_{j=1}^m\frac{\theta b_j\beta_j}{1-2\theta\lambda_j}+\frac12\sum_{j=1}^m\frac{\beta_j^2}{1-2\theta\lambda_j}\right]$.

This is the mgf of a multivariate normal with diagonal covariance and

$\mu_j(\theta)=\frac{\theta b_j}{1-2\theta\lambda_j},\qquad \sigma_j^2(\theta)=\frac1{1-2\theta\lambda_j}$.

\subsubsection*{Problem 9}

\term{Question.} Suppose $X_1,\ldots,X_n$ are independent $N(\mu,\sigma^2)$ observations and $\mu$ is known. (a) Show
$L(\sigma^2)\propto(\sigma^2)^{-n/2}\exp\left[-\frac1{2\sigma^2}\sum_{i=1}^n(X_i-\mu)^2\right].$

(b) If $\sigma^2\sim IG(\alpha,\beta)$, find the conditional distribution of $\sigma^2\mid X_1,\ldots,X_n$.

\term{Solution.} (a)

$L(\sigma^2)=\prod_{i=1}^n\frac1{\sqrt{2\pi}\sigma}\exp\left[-\frac{(X_i-\mu)^2}{2\sigma^2}\right]$

$\propto(\sigma^2)^{-n/2}\exp\left[-\frac1{2\sigma^2}\sum_{i=1}^n(X_i-\mu)^2\right].$

(b) The conditional density satisfies
$p(\sigma^2\mid x)\propto L(\sigma^2)p(\sigma^2)$

$\propto(\sigma^2)^{-n/2}\exp\left[-\frac1{2\sigma^2}\sum_{i=1}^n(X_i-\mu)^2\right](\sigma^2)^{-\alpha-1}e^{-\beta/\sigma^2}$

$\propto(\sigma^2)^{-(\alpha+n/2)-1}\exp\left[-\frac{\beta+\frac12\sum_{i=1}^n(X_i-\mu)^2}{\sigma^2}\right].$

Hence
$\sigma^2\mid X_1,\ldots,X_n\sim IG\left(\alpha+\frac n2,\beta+\frac12\sum_{i=1}^n(X_i-\mu)^2\right)$.

\subsubsection*{Problem 10}

\term{Question.} Given proposal transition matrix

$Q=\begin{pmatrix}1/6&0&1/2&1/3\\0&1/3&1/3&1/3\\0&1/2&0&1/2\\1/4&1/4&1/4&1/4\end{pmatrix}$,

use Metropolis-Hastings to construct a chain with limiting distribution $(1/6,1/6,1/3,1/3)$.

\term{Solution.} Let $Q=(q_{ij})$ and target probabilities

$\pi(1)=1/6,\quad\pi(2)=1/6,\quad\pi(3)=1/3,\quad\pi(4)=1/3$.

For $i\neq j$ and $q_{ij}\neq0$,
$\alpha_{ij}=\min\left[1,\frac{\pi(j)q_{ji}}{\pi(i)q_{ij}}\right].$

The nontrivial acceptance probabilities are

$\alpha_{13}=0,\quad\alpha_{14}=1,\quad\alpha_{23}=1,\quad\alpha_{24}=1,$

$\alpha_{32}=1/3,\quad\alpha_{34}=1/2,\quad\alpha_{41}=2/3,\quad\alpha_{42}=2/3,\quad\alpha_{43}=1$.

Thus

$\begin{array}{c|cccc}\alpha_{ij}&1&2&3&4\\\hline1&-&-&0&1\\2&-&-&1&1\\3&-&1/3&-&1/2\\4&2/3&2/3&1&-\end{array}$

and the revised probabilities are $p_{ij}=\alpha_{ij}q_{ij}$ for $i\neq j$, with $p_{ii}=1-\sum_{j\neq i}p_{ij}$. Hence
$P=\begin{pmatrix}2/3&0&0&1/3\\0&1/3&1/3&1/3\\0&1/6&7/12&1/4\\1/6&1/6&1/4&5/12\end{pmatrix}$.

This can be verified from $\pi P=\pi$.

\subsubsection*{Problem 11}

\term{Question.} In Example 12.6, sample a standard normal target

$f(x)=\frac1{\sqrt{2\pi}}e^{-x^2/2},\qquad -\infty<x<\infty,$

using Metropolis-Hastings with proposal

$q(x_n\mid x_{n-1})=\frac12e^{-|x_n|}$.

Starting from $X_0=0$, use proposal uniforms

\noindent\resizebox{\columnwidth}{!}{$\begin{array}{c|cccc}(U,V)&(0.18,0.87)&(0.02,0.42)&(0.31,0.08)&(0.12,0.79)\\W&0.55&0.03&0.18&0.43\end{array}$}

to generate four candidates, indicating acceptances and rejections. Here $U$ generates the exponential magnitude and $V$ determines its sign.

\term{Solution.} If $x_{k-1}$ is current and $x_k$ proposed, the Metropolis-Hastings ratio is

$\alpha=\frac{f(x_k)q(x_{k-1}\mid x_k)}{f(x_{k-1})q(x_k\mid x_{k-1})}$
$=\frac{e^{-x_k^2/2}e^{-|x_{k-1}|}}{e^{-x_{k-1}^2/2}e^{-|x_k|}}$.

The algorithm starts at $X_0=0$.

\noindent\resizebox{\columnwidth}{!}{$\begin{array}{c|cccc}k&1&2&3&4\\\hline(U,V)&(0.18,0.87)&(0.02,0.42)&(0.31,0.08)&(0.12,0.79)\\\text{candidate }x_k&-\log0.18=1.7148&\log0.02=-3.912&\log0.31=-1.1712&-\log0.12=2.1203\\\alpha&1.277&0.0186&1.2723&0.5418\\W&0.55&0.03&0.18&0.43\\\text{decision}&\text{accept}&\text{reject}&\text{accept}&\text{accept}\\X_k&1.7148&1.7148&-1.1712&2.1203\end{array}$}

The detailed ratios are

$k=1:\quad\alpha=\frac{e^{-(1.7148)^2/2}e^0}{e^0e^{-1.7148}}=1.277$;

$k=2:\quad\alpha=\frac{e^{-(-3.912)^2/2}e^{-1.7148}}{e^{-(1.7148)^2/2}e^{-3.912}}=0.0186$;

$k=3:\quad\alpha=\frac{e^{-(-1.1712)^2/2}e^{-1.7148}}{e^{-(1.7148)^2/2}e^{-1.1712}}=1.2723$;

$k=4:\quad\alpha=\frac{e^{-(2.1203)^2/2}e^{-1.1712}}{e^{-(-1.1712)^2/2}e^{-2.1203}}=0.5418$.
